\section*{Problemas}

\subsection*{Enunciados de los problemas}

% Recordar
\begin{problema}
	\label{prob:7ccb176}
	Enuncie, sin realizar ningún cálculo, la hipótesis nula de una prueba de homogeneidad para $c$ poblaciones independientes con una variable categórica de $r$ niveles, y la fórmula compacta del estadístico $\chi^2$ para el caso particular de $k$ proporciones ($r=2$).
\end{problema}

% Comprender
\begin{problema}
	\label{prob:44ef899}
	La fórmula matemática de cálculo ($\sum(O-E)^2/E$) es idéntica para las pruebas de independencia y de homogeneidad. Explique la diferencia conceptual profunda entre ambas en términos del \textbf{diseño de muestreo}: ¿qué está fijado de antemano por el investigador en cada caso (una sola muestra con dos variables observadas, o varias muestras independientes de tamaño prefijado)?
\end{problema}

% Aplicar
\begin{problema}
	\label{prob:3a99a23}
	En un ensayo clínico se contrasta una vacuna frente a un placebo tras seis meses de seguimiento: Grupo Vacuna ($n_1=150$): 12 infectados, 138 sanos. Grupo Placebo ($n_2=150$): 28 infectados, 122 sanos. Construya la tabla de contingencia $2\times2$ y ejecute la prueba $\chi^2$ de homogeneidad con $\nu=1$ para determinar si la tasa de infección difiere significativamente entre ambos grupos al $\alpha=0.05$ (use $\chi^2_{0.05,1}=3.841$).
\end{problema}

% Analizar
\begin{problema}
	\label{prob:a98f897}
	Con los datos del problema anterior (vacuna: 12/150 infectados; placebo: 28/150 infectados), calcule el estadístico normal tipificado $Z$ de una prueba bilateral convencional de dos proporciones (con error estándar agrupado), y verifique numéricamente la identidad algebraica $Z^2\equiv\chi^2$ frente al estadístico $\chi^2$ calculado en el problema anterior. Explique brevemente por qué esta identidad se cumple siempre en tablas $2\times2$.
\end{problema}

% Evaluar
\begin{problema}
	\label{prob:cb7e8b6}
	Un analista, tras rechazar la hipótesis global $H_0:p_1=p_2=p_3=p_4$ en una prueba de homogeneidad de $k=4$ proporciones, decide identificar cuáles proporciones difieren realizando las $\binom{4}{2}=6$ pruebas $Z$ bilaterales posibles entre pares, cada una al nivel individual $\alpha=0.05$, sin ningún ajuste adicional. Evalúe esta práctica: ¿qué problema estadístico introduce, y por qué el procedimiento de comparación múltiple de Marascuilo (que usa el rango crítico $r_{ij}=\sqrt{\chi^2_{\alpha,k-1}}\sqrt{\hat p_i(1-\hat p_i)/n_i+\hat p_j(1-\hat p_j)/n_j}$) es la alternativa estadísticamente correcta?
\end{problema}

% Crear
\begin{problema}
	\label{prob:4eeb5bb}
	Diseñe un ejemplo original con $k=3$ poblaciones independientes (contexto y proporciones muestrales de éxito) donde se aplique la prueba de homogeneidad de $k$ proporciones, distinto del ejemplo de vacunación ya visto en esta sección. Plantee $H_0$ y calcule la proporción global combinada $\bar p$.
\end{problema}

\subsection*{Soluciones detalladas}

% Recordar
\begin{solproblema}[prob:7ccb176]
	$H_0: p_{i1}=p_{i2}=\cdots=p_{ic}$ para cada categoría $i=1,\ldots,r$ (las $c$ poblaciones tienen idéntica distribución categórica). Para $k$ proporciones ($r=2$): $\chi^2 = \sum_{j=1}^k \dfrac{(X_j-n_j\bar p)^2}{n_j\bar p(1-\bar p)}$, con $\bar p=\sum X_j/\sum n_j$.
\end{solproblema}

% Comprender
\begin{solproblema}[prob:44ef899]
	En la prueba de \textbf{independencia}, se extrae \textbf{una sola muestra} de tamaño $N$ y se observan \textbf{dos variables categóricas simultáneamente} sobre cada individuo; ni los totales de fila ni los de columna están fijados de antemano, ambos emergen de la muestra observada. En la prueba de \textbf{homogeneidad}, en cambio, se extraen de antemano $c$ muestras \textbf{independientes} de tamaños prefijados por el diseño experimental ($n_1,\ldots,n_c$, es decir, los totales de columna están fijos \emph{a priori}, no observados), y se mide una única variable categórica dentro de cada una. La fórmula de cálculo coincide porque, condicionalmente a los totales marginales observados, ambos escenarios producen la misma distribución muestral asintótica del estadístico — pero el diseño de muestreo subyacente (y por tanto la pregunta de investigación que responde la prueba) es distinto.
\end{solproblema}

% Aplicar
\begin{solproblema}[prob:3a99a23]
	\begin{center}
		\begin{tabular}{|l|c|c|c|}\hline
			 & Infectados & Sanos & Total \\\hline
			Vacuna & 12 & 138 & 150 \\\hline
			Placebo & 28 & 122 & 150 \\\hline
			Total & 40 & 260 & 300 \\\hline
		\end{tabular}
	\end{center}
	$E_{11}=E_{21}=150(40)/300=20$; $E_{12}=E_{22}=150(260)/300=130$.
	\begin{align*}
		\chi^2 = \frac{(12-20)^2}{20}+\frac{(138-130)^2}{130}+\frac{(28-20)^2}{20}+\frac{(122-130)^2}{130} = 3.20+0.492+3.20+0.492 = 7.385.
	\end{align*}
	Con $\chi^2_{0.05,1}=3.841$: como $7.385>3.841$, \textbf{se rechaza} $H_0$: la vacuna reduce significativamente la tasa de infección.
\end{solproblema}

% Analizar
\begin{solproblema}[prob:a98f897]
	$\hat p_1=12/150=0.08$, $\hat p_2=28/150\approx0.18667$. Proporción combinada $\hat p=40/300\approx0.13333$. Error estándar agrupado:
	\begin{align*}
		\text{SE}_p = \sqrt{\hat p(1-\hat p)\left(\frac{1}{150}+\frac{1}{150}\right)} \approx0.03925.
	\end{align*}
	$Z = (0.08-0.18667)/0.03925 \approx-2.7175$, así $Z^2\approx7.385$, coincidiendo (salvo redondeo) con $\chi^2=7.385$. Esta identidad se cumple siempre en tablas $2\times2$ porque, por definición, el cuadrado de una Normal estándar sigue una $\chi^2_1$ ($Z\sim N(0,1)\implies Z^2\sim\chi^2_1$), y el estadístico de Pearson en una tabla $2\times2$ es estructuralmente la versión al cuadrado de la brecha estandarizada entre las dos proporciones binomiales.
\end{solproblema}

% Evaluar
\begin{solproblema}[prob:cb7e8b6]
	La práctica es \textbf{estadísticamente incorrecta}. Realizar $6$ pruebas independientes cada una al nivel $\alpha=0.05$ infla la Tasa de Error Familiar (FWER): la probabilidad de cometer \emph{al menos un} falso positivo entre las $6$ comparaciones es considerablemente mayor al $5\%$ nominal (similar al problema de comparaciones múltiples visto para pruebas de hipótesis independientes). El procedimiento de \textbf{Marascuilo} corrige esto usando el cuantil global $\chi^2_{\alpha,k-1}$ (no un cuantil individual por par) dentro del rango crítico $r_{ij}$, garantizando que la probabilidad conjunta de declarar \emph{cualquier} par como significativamente distinto sin serlo realmente se mantenga acotada en el nivel $\alpha$ global — es la alternativa correcta porque controla el error a nivel del conjunto completo de comparaciones, no prueba por prueba.
\end{solproblema}

% Crear
\begin{solproblema}[prob:4eeb5bb]
	\textbf{Ejemplo:} un estudio de retención de clientes compara la tasa de cancelación de suscripción en $k=3$ planes de precio: Plan Básico ($n_1=300$, $X_1=45$ cancelaciones), Plan Estándar ($n_2=350$, $X_2=28$), Plan Premium ($n_3=200$, $X_3=10$). $H_0: p_1=p_2=p_3$ (la tasa de cancelación es la misma en los tres planes). La proporción global combinada es:
	\begin{align*}
		\bar p = \frac{45+28+10}{300+350+200} = \frac{83}{850} \approx0.0976.
	\end{align*}
\end{solproblema}
